% ======================================================================
% Conceptual companion to the manuscript: the physics behind eddy-current
% strain reconstruction in LPBF, explained with text and figures.
% Build:  latexmk -pdf concepts.tex   (figures must be compiled first; `make`)
% ======================================================================
\documentclass[11pt,a4paper]{article}

\usepackage[margin=2.6cm]{geometry}
\usepackage{amsmath,amssymb,bm}
\usepackage{graphicx}
\usepackage{siunitx}
\usepackage{xcolor}
\usepackage[hidelinks]{hyperref}
\usepackage[capitalise]{cleveref}

\graphicspath{{figures/}}

\newcommand{\dsig}{\Delta\sigma}
\newcommand{\eps}{\bm{\varepsilon}}
\newcommand{\sig}{\bm{\sigma}}

\title{\textbf{The physics of eddy-current strain sensing\\
in laser powder bed fusion}\\[4pt]
\large A conceptual companion to the research article}
\author{M.~Carraturo \and B.\,I.~Ates}
\date{\today}

\begin{document}
\maketitle

\noindent
This note explains, concept by concept, the physics that the companion
research article assembles into a measurement chain: how a coil senses the
electrical conductivity of a metal part, how mechanical strain alters that
conductivity, why a single measurement can never determine a full strain
tensor, and how missing information is supplied by a process simulation.
\Cref{fig:overview} shows where each concept sits in the chain.

\begin{figure}[htb]
  \centering
  \includegraphics{fig_overview}
  \caption{The measurement chain. A rosette of directional eddy-current
  probes rides on the recoater of a laser powder bed fusion (LPBF) machine
  and scans every layer; the readings are fused with a thermomechanical
  process simulation to reconstruct the strain-tensor field
  $\varepsilon_{ij}(\bm x)$. Sections~\ref{sec:em}--\ref{sec:lpbf} unpack
  the physics behind each link.}
  \label{fig:overview}
\end{figure}

% ======================================================================
\section{Electromagnetic foundations}\label{sec:em}
% ======================================================================

\subsection{Eddy currents}
A coil carrying an alternating current $I\,e^{\mathrm{j}\omega t}$ produces an
alternating magnetic field $\bm B(t)$. Faraday's law,
\begin{equation}
  \nabla\times\bm E = -\frac{\partial \bm B}{\partial t},
  \label{eq:faraday}
\end{equation}
states that a time-varying magnetic field induces an electric field; in a
conductor of conductivity $\sigma$ that electric field drives circulating
currents, $\bm J_{\mathrm e}=\sigma\bm E$ --- the \emph{eddy currents}
(\cref{fig:eddy}). By Lenz's law their own magnetic field opposes the change
of flux that created them. Seen from the coil's terminals this has two
effects: the opposing flux reduces the coil's effective inductance, and the
Joule dissipation $\sigma|\bm E|^2$ of the eddy currents appears as an added
resistance. The conductor therefore leaves a fingerprint on the coil
\emph{impedance}, and that fingerprint encodes $\sigma$ --- no contact
required.

\begin{figure}[htb]
  \centering
  \includegraphics{fig_eddy_currents}
  \caption{Eddy-current induction, drawn in the meridional plane of an
  axisymmetric coil. The alternating drive current (cross: into the page;
  dot: out of the page) produces the primary field $\bm B(t)$; the induced
  eddy currents $\bm J_{\mathrm e}$ circulate in the opposite sense and are
  confined to a thin layer of thickness ${\sim}\delta$ below the surface.}
  \label{fig:eddy}
\end{figure}

\subsection{The skin effect}
Inside the conductor, the field does not penetrate freely: the induced
currents screen it, and the governing equations reduce to a diffusion
problem whose solution decays exponentially with depth $z$,
\begin{equation}
  |J|(z) = J_0\, e^{z/\delta},
  \qquad
  \delta=\sqrt{\frac{2}{\omega\mu\sigma}},
  \label{eq:skin}
\end{equation}
with permeability $\mu$ ($=\mu_0$ for non-magnetic 316L steel). The
\emph{skin depth} $\delta$ is the single most important length scale of the
method: it is where the measurement lives. \Cref{fig:skin} plots
\cref{eq:skin} for 316L ($\sigma\approx\SI{1.35}{\mega\siemens\per\metre}$):
$\delta$ shrinks from \SI{1.8}{\milli\metre} at \SI{60}{\kilo\hertz} to
\SI{0.4}{\milli\metre} at \SI{1}{\mega\hertz}. Two practical consequences
follow. First, the excitation frequency is a \emph{depth dial}: low
frequencies look deep, high frequencies hug the surface, which the
finite-element model of the article confirms (\cref{fig:depth}). Second,
even the thinnest practical skin depth spans roughly ten LPBF layers
($30$--$\SI{50}{\micro\metre}$ each), so a single reading is always an
average over many layers.

\begin{figure}[htb]
  \centering
  \includegraphics{fig_skin_effect}
  \caption{The skin effect in 316L: current density versus depth,
  \cref{eq:skin}, at three excitation frequencies. At the depth $z=-\delta$
  the current density has fallen to $e^{-1}\approx0.37$ of its surface
  value.}
  \label{fig:skin}
\end{figure}

\begin{figure}[htb]
  \centering
  \includegraphics{fig_depth}
  \caption{The same depth selectivity computed with the finite-element
  forward model of the companion article: mean sensitivity versus depth at
  two frequencies. Frequency is a design variable that chooses how deep the
  probe looks.}
  \label{fig:depth}
\end{figure}

\subsection{The impedance plane}
The coil's impedance $Z=R+\mathrm{j}X$ combines the resistance $R$ (loss)
and the reactance $X=\omega L$ (magnetic energy storage). In air the coil
is almost purely inductive, $Z_0\approx\mathrm{j}\omega L_0$. Bringing a
conductor underneath moves the operating point as sketched in
\cref{fig:zplane}: the eddy-current loss adds resistance
($\Delta R>0$) while the opposing flux removes reactance ($\Delta X<0$). As
the conductivity (or the frequency) grows, the point traces a
characteristic arc in the impedance plane. That arc is the heart of the
inverse problem: measuring where $Z$ sits on it tells us $\sigma$. To
compare different coils and instruments the article works with the
normalised impedance
\begin{equation}
  Z_{\mathrm{norm}}
  = \frac{Z_0 - Z}{\operatorname{Im} Z_0},
  \label{eq:znorm}
\end{equation}
a dimensionless quantity insensitive to the number of turns and to overall
scale factors.

\begin{figure}[htb]
  \centering
  \includegraphics{fig_impedance_plane}
  \caption{The impedance plane. A conductor moves the coil's operating
  point from the air value $Z_0$ (purely inductive) to $Z$: the eddy-current
  loss adds resistance, the flux opposition removes reactance. Increasing
  conductivity or frequency drives the point along a characteristic locus
  --- inverting that locus recovers $\sigma$ from the measured $Z$.}
  \label{fig:zplane}
\end{figure}

\subsection{What exactly does the probe see?}
A perturbation argument makes the sensing precise. If the conductivity
changes by a small amount $\delta\sigma(\bm x)$ somewhere in the part, the
impedance changes by
\begin{equation}
  \delta Z = -\frac{1}{I^2}\int \delta\sigma\;
  \bm E\!\cdot\!\bm E \,\mathrm{d}V :
  \label{eq:recip}
\end{equation}
the probe weighs every point of the part by the \emph{square of the local
electric field} it drives there. For a circular coil over a flat surface
the driven field is purely tangential (azimuthal), with two consequences
visible in \cref{fig:kernel}. Spatially, the sensitivity is concentrated
under the winding and within one skin depth of the surface.
Directionally --- and this is the crucial point for strain sensing --- the
kernel $\bm E\!\cdot\!\bm E$ contains only \emph{in-plane} field components:
a conventional coil is structurally blind to the through-thickness
conductivity $\sigma_{zz}$, no matter how sensitive its electronics.

\begin{figure}[htb]
  \centering
  \includegraphics{fig_kernel}
  \caption{Sensitivity density $|E_\phi|^2$ of a circular coil, computed
  with the finite-element model: the probe ``sees'' a ring under the
  winding, one skin depth deep. Because the field is purely in-plane, the
  probe cannot sense the out-of-plane conductivity component.}
  \label{fig:kernel}
\end{figure}

% ======================================================================
\section{From strain to conductivity --- and back}\label{sec:strain}
% ======================================================================

\subsection{Why strain changes conductivity}
The resistance of a conducting bar is $R=\rho L/A$, with resistivity
$\rho=1/\sigma$, length $L$ and cross-section $A$. Stretching the bar by a
strain $\varepsilon$ lengthens the current path, $L\to L(1+\varepsilon)$,
and thins it, $A\to A(1-\nu\varepsilon)^2$ with Poisson ratio $\nu$
(\cref{fig:bar}). On top of this purely geometric effect, the resistivity
itself changes: elastic strain modifies the lattice spacing and with it the
scattering of conduction electrons. The total effect --- geometry plus
intrinsic piezoresistivity --- is the working principle of the strain
gauge, and it is linear for small strains. Expressed for the conductivity,
this is the scalar law used in the article,
\begin{equation}
  \frac{\dsig}{\sigma_0} = \kappa\,\varepsilon,
  \qquad \kappa = -0.326 \;\;\text{(316L, uniaxial calibration)},
  \label{eq:scalar}
\end{equation}
where $\sigma_0$ is the unstrained conductivity: tensile strain lowers the
conductivity by about a third of the strain value.

\begin{figure}[htb]
  \centering
  \includegraphics{fig_piezoresistive_bar}
  \caption{Why strain changes resistance. Stretching lengthens and thins
  the current path (geometric effect) and alters the resistivity itself
  (intrinsic piezoresistive effect); both are linear in $\varepsilon$ for
  elastic strains.}
  \label{fig:bar}
\end{figure}

\subsection{Strain is a tensor --- so conductivity becomes one}
A part in three dimensions is not a bar. Its strain state is a symmetric
tensor $\eps$ with six independent components: three normal strains and
three shears. A material stretched along $x$ conducts differently along
$x$ than along $y$ --- strain \emph{breaks the isotropy} of the material,
and the conductivity must be described by a symmetric tensor
$\sigma_{ij}$ as well. For an isotropic metal the linear coupling between
the two tensors has exactly two independent constants: a longitudinal
coefficient $\kappa_\parallel$ (response along the strained direction) and
a transverse coefficient $\kappa_\perp$ (response perpendicular to it).
The scalar $\kappa$ of \cref{eq:scalar} is then revealed to be a
\emph{projection}: in a uniaxial test, the conductivity change measured
along the load is $\kappa_\parallel-2\nu\kappa_\perp$, while a probe
oriented transversely reads a different --- even opposite-signed ---
value. \Cref{fig:rosette_curve} shows this ``conductivity rosette'': the
same strain state produces an effective $\kappa$ anywhere between $-0.35$
and $+0.06$ depending on the sensing direction. A single scalar reading is
therefore fundamentally ambiguous about the underlying strain tensor.

\begin{figure}[htb]
  \centering
  \includegraphics{fig_rosette}
  \caption{The conductivity rosette: effective $\kappa$ sensed in a
  uniaxial-stress test as a function of the in-plane sensing angle. One
  number cannot summarise this curve --- which is precisely why one probe
  cannot determine a strain tensor.}
  \label{fig:rosette_curve}
\end{figure}

\subsection{Directional probes: a rosette for conductivity}
Mechanical engineers solved the analogous problem a century ago: a single
strain gauge reads one direction, so three gauges at $0^\circ/60^\circ/
120^\circ$ --- a \emph{rosette} --- determine the full in-plane strain
state. The same idea transfers to conductivity (\cref{fig:probes}). A
directionally selective probe driving current mainly along the unit vector
$\hat{\bm n}$ reads the projection
\begin{equation}
  y(\hat{\bm n})=\hat{\bm n}^{\!\top}\sig\,\hat{\bm n},
  \label{eq:proj}
\end{equation}
and a set of differently oriented probes assembles a linear system for the
components of $\sig$ --- and, through the elastoresistive coupling, for
$\eps$. Linear algebra then answers ``how many components can this probe
set resolve?'' exactly: the rank of the measurement map counts the
resolvable degrees of freedom, and its condition number measures how
strongly noise is amplified. Three in-plane probes resolve the three
in-plane components; tilted or side-mounted probes are needed for the
remaining three, and from the top surface alone those remain weakly
observable --- the skin-effect physics of \cref{sec:em} is what removes
them.

\begin{figure}[htb]
  \centering
  \includegraphics{fig_probe_rosette}
  \caption{A conductivity rosette: three directional probes at
  $0^\circ/60^\circ/120^\circ$, each reading the projection
  $y_k=\hat{\bm n}_k^{\!\top}\sig\,\hat{\bm n}_k$ of the conductivity
  tensor. The arrangement mirrors the classical strain-gauge rosette.}
  \label{fig:probes}
\end{figure}

\subsection{Filling the blind spots: fusion with a simulation prior}
What the probes cannot see must come from somewhere else. LPBF is unusual
among manufacturing processes in that a physics-based prediction of the
strain field --- a thermomechanical \emph{process simulation} --- is
routinely available. Bayesian estimation provides the principled way to
combine the two information sources: the simulation supplies a \emph{prior}
probability for the strain, the ECT readings supply a \emph{likelihood},
and the product of the two is the \emph{posterior} --- the best estimate
given everything known (\cref{fig:bayes}). For Gaussian uncertainties the
posterior has a closed form, the maximum-a-posteriori estimate of the
article, and an attendant \emph{resolution matrix} that states, component
by component, what fraction of the answer came from the data and what
fraction from the prior. Well-observed components (the in-plane strains)
are corrected by the measurements even when the simulation is biased;
blind components (out-of-plane) gracefully fall back to the simulation.

\begin{figure}[htb]
  \centering
  \includegraphics{fig_bayes_fusion}
  \caption{Data assimilation in one dimension. The process simulation
  provides a broad prior; the ECT measurement provides a narrower
  likelihood; the posterior combines them, weighted by their confidences.
  Where the data are uninformative the posterior simply returns the prior
  --- the mechanism by which unobservable strain components are filled in.}
  \label{fig:bayes}
\end{figure}

% ======================================================================
\section{The LPBF context}\label{sec:lpbf}
% ======================================================================

\subsection{Where the residual strain comes from}
LPBF melts metal powder layer by layer with a laser. Each pass heats a
shallow zone to extreme temperatures within milliseconds. The hot material
wants to expand but is constrained by the cold, stiff material beneath: it
yields in compression (\cref{fig:tgm}a). When it cools it wants to shrink
below its original length --- and is again constrained, so the contraction
is locked in as \emph{tensile} residual stress near the surface, balanced
by compression below (\cref{fig:tgm}b). Repeated over thousands of layers,
this temperature-gradient mechanism builds residual elastic strains of
order $10^{-3}$ throughout the part --- large enough to distort or crack
it, and exactly the quantity the measurement chain of
\cref{fig:overview} is designed to map.

\begin{figure}[htb]
  \centering
  \includegraphics{fig_tgm}
  \caption{The temperature-gradient mechanism that generates residual
  strain in LPBF. (a) The laser-heated layer cannot expand freely and
  yields in compression. (b) On cooling, the hindered contraction remains
  as tension near the surface, balanced by compression underneath.}
  \label{fig:tgm}
\end{figure}

\subsection{Why the measurement is hard --- by the numbers}
The concepts above also expose the difficulty. Residual strains of
$1$--$2\times10^{-3}$, multiplied by $|\kappa|\approx0.33$, give relative
conductivity changes of only $3$--$7\times10^{-4}$ --- parts in ten
thousand. Meanwhile the conductivity of 316L drifts with temperature at
roughly $10^{-3}$ per kelvin: a \SI{1}{\kelvin} drift mimics the entire
strain signal. Plastic deformation and dislocation density shift the
conductivity too, entangling the elastic-strain signal with
microstructure. And the skin depth of \cref{eq:skin} means each reading
averages tens of layers in depth. None of these is a show-stopper; together
they specify the instrument: a \emph{differential, thermally referenced},
\emph{multi-frequency}, \emph{directional} probe rosette --- which is
precisely the configuration whose observability and inversion the
companion article quantifies.

\end{document}
